% 2009 年湖南卷文科第 7 题：四个候选的 y=f(x) 图象（A）—（D），选项本身即图。
% 来源 PDF：普通高考/2009/2009湖南文.pdf，SHA-256 db41a34b5d57a8db3f2ee49872e3dbc1e290c2e7390fc96cf4e084b5a679e3eb。
% 原页：第 1 页（A3 三栏）第 1 栏下部，书页码 529。
% 重绘依据：四图均按原页逐幅核对凹凸与端点——
%   (A) 自 a 到 b 始终下凸（斜率递增），起端较平、末端陡；
%   (B) 自 a 到 b 始终上凸（斜率递减），起端陡、末端渐平；
%   (C) 为一条直线段，斜率为常数；
%   (D) 呈 S 形：先缓、中段陡、末段又渐平，变号点使 f' 先增后减。
%   四图均在 x=a、x=b 处由曲线端点向 x 轴作虚线，O 为原点，箭头指向正方向。
\begin{tikzpicture}[semithick, every node/.style={inner sep=1.3pt, font=\small}]

  %=========================== (A) 下凸 ===========================
  \begin{scope}[shift={(0,0)}]
    \draw[-{Stealth}] (-0.40,0) -- (4.45,0) node[right] {$x$};
    \draw[-{Stealth}] (0,-0.35) -- (0,2.75) node[above] {$y$};
    \node[below left=1pt] at (0,0) {$O$};
    \draw[dashed] (0.90,0) -- (0.90,0.55);
    \draw[dashed] (3.30,0) -- (3.30,2.28);
    \draw (0.90,0.55) .. controls (1.85,0.62) and (2.60,1.15) .. (3.30,2.28);
    \node[below] at (0.90,0) {$a$};
    \node[below] at (3.30,0) {$b$};
    \node at (-0.45,-0.55) {(A)};
  \end{scope}

  %=========================== (B) 上凸 ===========================
  \begin{scope}[shift={(5.60,0)}]
    \draw[-{Stealth}] (-0.40,0) -- (4.45,0) node[right] {$x$};
    \draw[-{Stealth}] (0,-0.35) -- (0,2.75) node[above] {$y$};
    \node[below left=1pt] at (0,0) {$O$};
    \draw[dashed] (0.90,0) -- (0.90,0.35);
    \draw[dashed] (3.30,0) -- (3.30,2.20);
    \draw (0.90,0.35) .. controls (1.35,1.60) and (2.25,2.02) .. (3.30,2.20);
    \node[below] at (0.90,0) {$a$};
    \node[below] at (3.30,0) {$b$};
    \node at (-0.45,-0.55) {(B)};
  \end{scope}

  %=========================== (C) 直线段 ===========================
  \begin{scope}[shift={(0,-4.55)}]
    \draw[-{Stealth}] (-0.40,0) -- (4.45,0) node[right] {$x$};
    \draw[-{Stealth}] (0,-0.35) -- (0,2.75) node[above] {$y$};
    \node[below left=1pt] at (0,0) {$O$};
    \draw[dashed] (0.90,0) -- (0.90,0.45);
    \draw[dashed] (3.30,0) -- (3.30,2.20);
    \draw (0.90,0.45) -- (3.30,2.20);
    \node[below] at (0.90,0) {$a$};
    \node[below] at (3.30,0) {$b$};
    \node at (-0.45,-0.55) {(C)};
  \end{scope}

  %=========================== (D) S 形 ===========================
  \begin{scope}[shift={(5.60,-4.55)}]
    \draw[-{Stealth}] (-0.40,0) -- (4.45,0) node[right] {$x$};
    \draw[-{Stealth}] (0,-0.35) -- (0,2.75) node[above] {$y$};
    \node[below left=1pt] at (0,0) {$O$};
    \draw[dashed] (0.90,0) -- (0.90,0.55);
    \draw[dashed] (3.30,0) -- (3.30,2.35);
    \draw (0.90,0.55) .. controls (1.65,0.72) and (2.30,2.10) .. (3.30,2.35);
    \node[below] at (0.90,0) {$a$};
    \node[below] at (3.30,0) {$b$};
    \node at (-0.45,-0.55) {(D)};
  \end{scope}
\end{tikzpicture}
